\section{\pysv\ 1.2}
\setlength{\footmarkwidth}{1.8em}
\pagewiselinenumbers
\newlength{\dawidth}
\settowidth{\dawidth}{ड}
\addtolength{\dawidth}{0.08em}
\XeTeXglyph 101 \raisebox{0.25ex}{\scalebox{1}[0.8]{ड}}\kern -\dawidth \raisebox{-0.65ex}{\scalebox{1}[0.8]{भ}} \mll{ONETWOONE}\mula{सर्व्ववृत्ति\vrtms{निरोधे}{TM\Me}{\rdg{॰निरोधे प्य}{L}} त्वसम्प्रज्ञात इति}\donote{ONETWOONE}{{\dn सर्व्ववृत्तिनिरोधे त्वसम्प्रज्ञात इति}}।\llbl{SUTR2INTRO} तुशब्दो ऽवधारणार्त्थः। \mll{SUTR2INTRO}\mula{तस्यै}वं\vrtms{विधस्य}{\Me(em.)}{\rdg{॰वंविषयस्य}{TLMA}} समाधेरसम्प्रज्ञातस्यैव केवलस्य \mula{लक्षणाभिधित्सयेदं सूत्रम्प्रववृते}\donote{SUTR2INTRO}{{\dn तस्य लक्षणाभिधित्सयेदं सूत्रम्प्रववृते}} प्रवृत्तम्---
% evaṃviṣayasya is in all the manuscripts
% I consider "eva" to be part of the mūla. "kevalasya" appears a gloss for "eva". Otherwise "tasyaiva kevalasya" appears slightly redundant.
\newfontfamily
\str{योगश्चित्तवृत्तिनि\mms रोध (२)} \mds \vrt{इति}{L}{\rdg{ईति वक्तव्यं}{M}, \rdg{ई[[ती]]ति वक्तव्यम्}{\Me}}॥

ननु च त\mte स्य लक्षणमभिधातुमेतत्सूत्रम्प्रवृत्तम? तथा च सति योगस्य चित्तवृत्तिनिरोध इति वक्तव्यम्। 
सामानाधि\vrtms{करण्यन्न}{LM\Me\vv{॰ण्यं न}{\Me}}{\rdg{॰करण्यान्न}{\Td}} युक्तम्। न हि लक्ष्यमेव लक्षणं स्यात्। चित्तवृत्तिनिरोधलक्षण इति वा वक्तव्यम्॥
% na hi lakṣyam eva lakṣaṇaṃ syāt | cittavṛtti... is awkward.

नैष दोषः। लक्ष्ये ल\mms क्षण\mds ाध्यासात्। यथा\mte यम्पिण्डो देवदत्त इति लक्षये लक्षणम्प्रत्यस्यते॥
% piṇḍa and Devadatta: BAUBh, MBh, Āryabhaṭīya

ननु \mms च स\mme र्व्ववृत्तिनिरोधे त्वसम्प्रज्ञातः। तस्य चे\mms दं लक्षणम्। ततश्च योगः सर्व्वचि\mme त्तवृत्तिनिरोध इति वक्तव्यम। न चोक्तम्। अतो न व्या\mds पि \mms लक्षणम्॥

सर्व्वशब्दाग्रहणे तावत्कारणमुच्यते---यदि सर्व\mte ग्रहणं क्रियेत तदैकदेशनिरोधनिमित्तस्य सम्प्रज्ञातस्य योगत्वन्निवर्तितं स्यात्। तन्मा भूदिति सर्व्वशब्दाग्रहणम्॥

अतस्तर्हि सम्प्रज्ञातासम्प्रज्ञातयोः सामान्यलक्षणम्, विशेषानुपादानात्॥

नैष दोषः, योगे प्रकृते पुनर्य्योगग्रहणात्। वि\mts शेषाविवक्षायां हि यो\mde ग\mme ग्रहणन्नार्त्थवत्स्यात्॥

ननु चानुशासने गुणभूतत्वाद्योगस्य, इह योगग्रहणे ऽसति चित्तवृत्तिनिरोध इत्युक्ते ऽनुशासनार्त्थता भवेत्सूत्रस्य॥

नैवम्, योगानुशासनस्य प्रस्तुतत्वात्। योगानुशासनं हि प्रारब्धन्नायो\mts गानुशासनम्। त\mde ल्ल\mme क्ष\-णाभिधाने 
ह्यतिप्रसंगः पदार्त्थासंगतिश्च। न हि चित्तवृत्तिनिरोधो योगानुशासनमिति पदार्त्थः संगच्छते॥

एवमपि योगग्रहणात्कुतो ऽसम्प्रज्ञातस्यैव लक्षणम्, न पुनः सम्प्रज्ञातस्य लक्षणं स्यात्? योगग्रहणस्याविशिस्ष्टत्वादित्युच्यते---

न, व्यभिचारात्। [यद्]यद्व्यभिचरति न तत्तस्य लक्षणम्। निरोधस्त्वसम्प्रज्ञाते ऽपि वर्त्तत इति सम्प्रज्ञातं व्यभिचरति। न 
हि विषाणित्वं गोर्लिंगम्, महिष्यादिभिर्व्यभिचारात्॥

ननु चासम्प्रज्ञातमपि व्यभिचरति, सम्प्रज्ञातस्यापि हेतुत्वात्?

सत्यमेवम्। किन्तु न निरोधादन्येन लक्षयितुमसम्प्रज्ञातः शक्यते। निरोध एव तु तस्य लक्षणम्, अन्याभावात्, सम्प्रज्ञातस्य त्वसाधारणवितर्कादिलक्षणलक्ष्यत्वात्। यथा `स्पर्शलक्षणमि'त्युक्ते वायोरेवासाधारण्याल्लक्ष्यते स्पर्शवत्त्वेन। अथ च विद्यत एव स्पर्शवत्वन्तेजःप्रभृतिष्वपि। एवञ्च सति योगग्रहणमनुवादमात्रम्। अत एव च सर्वग्रहणन्न कृतम्। सूत्रे निरोधग्रहणेनैवासम्प्रज्ञातलक्षणत्वे सिद्धे यदि सर्व्वग्रहणं क्रियेत तत्सम्प्रज्ञातस्य योगतां निवारयेत्। तस्मादसम्प्रज्ञातस्यैवेदं लक्षणम्॥

तथा चोपसंहरति भाष्यकारः---न तत्र किञ्चित्सम्प्रज्ञायत इत्यसम्प्रज्ञातः। स योगश्चित्तवृत्तिनिरोध इति। 
चित्ततद्वृत्तितन्निरोधव्याचिख्यासया चित्तं हीत्यादि भाष्यम्। किमात्मकम्पुनस्तच्चित्तं यस्य वृत्तिनिरोधो योग इत्यत आह---
चित्तं हि प्रख्याप्रवृ\linelabel{tmrepeats}त्तिस्थितिशीलत्वात्त्रिगुणमिति॥

तत्र `चित्तन्त्रिगुणमि'ति प्रतिज्ञायते। `प्रख्याप्रवृत्तिस्थितिशीलत्वादि'ति हेतुः। प्रख्या प्रख्यानम्प्रकाशनम्। सा(?) हि 
सत्वगुणस्य धर्म्मः। प्रवृत्तिः प्रवर्त्तनं व्यापारः। स हि रजोधर्म्मः। स्थितिः स्थानं वरणम्प्रतिबन्ध इति। स च तमोधर्म्मः। 
एवंशीला सत्वादयो गुणाः। चित्तञ्च प्रख्यादिशीलम्। तस्मात्त्रयाणां गुणानाम् परिणामश्चित्तम्भवितुमर्हति॥ % NEED A NOTE ON SAA OR SA, REFERENCES TO THOSE DHARMAS OF GU.NAS!

तथा प्रख्यादिहेतोरसिद्धतामाशंक्याह---प्रख्यास्वरूपं हि चित्तसत्वमिति। हिशब्दः प्रसिद्धावद्योतनार्त्थः। प्रसिद्धं हि लोके 
शास्त्रे च सर्व्वावभासकत्वञ्चित्तस्य। चित्तमेव सत्वञ्चित्तसत्वम्, सत्वप्रधानगुणपरिणामत्वात्॥ % need a big note on cittasya sarvaavabhaasakatvam!

अथ वान्यसंसर्गवियोगाभ्यामनेकवृत्तिख्यातिमात्रत्वञ्चित्तस्य स्पष्टन्दर्शयितुं शक्यत इति ब्रवीति---प्रख्यास्वरूपं हि 
चित्तसत्वमिति॥

इदानीं वृत्तयो व्याख्यायन्ते। तासाञ्चानेकत्वं रजस्तमसोरुद्भवाभिभवनिमित्तं विरुद्धत्वञ्चेत्याह---रजस्तमोभ्यां संसृष्टं 
सम्प्रधानाभ्यां यदा संसृष्टं तदानीमैश्वर्यविषयप्रियम्भवति। प्रीतिः रागः। ऐश्वर्यविषयरागयुक्ता वृत्तयो भवन्तीत्यर्त्थः। 
तत्तम्सानुविद्धम्। तदेव प्रख्यास्वरूपञ्चित्तं यदा गुणभूतरजस्केनोद्भूतेन तमसानुविद्धम्, तदानीमधर्मादिचतुष्टयविषया वृत्तयः क्
लिष्टा जायन्ते॥

तदेव प्रक्षीणमोहावरणामिति। न्यग्भूततमस्कन्तमःक्षयादेव जलधारापगमादिव रविबिम्बं सर्वतः 
प्रद्योतमानं रजोमात्रया रजोलेशेनानुविद्धन्तदा धर्म्माद्युपगतं, धर्म्मादिचतुष्टयविषया वृत्तयो ऽक्लिष्टा भवन्ति॥

यदा तदेव रजोमलापेतं स्वरूपप्रतिष्ठं केवलेन स्वेन प्रख्यारूपेणावस्थितन्तदा सत्वपुरुषान्यताख्यातिमात्रम्भवति। सत्वञ्चित्तम्, 
पुरुषो भक्ता, तयोरन्यथाविविक्तता, तत्ख्यातिस्तदवगमः। तन्मात्रग्रहणं \kle शाद्यभावप्रदर्शनार्त्थम्। धर्म्ममेघध्यानोपगतम्भवति। 
धर्म्ममेघो नाम समाधिः। तदेवम्प्रसंख्यानबलाद्रजस्तमसी तिरस्कृत्य केवलेन ख्यात्यात्मना पुरुषस्वरूपमात्रदर्शनेनावस्थानध्यानं 
प्रसंख्यानमित्याचक्षते ध्यायिनो योगिनः॥

एवं वृत्तिस्वरूपे व्याख्याते तन्निरोधप्रदर्शनार्त्थमाह---चितिशक्तिरित्यादि। चितिरेव शक्तिः चितिशक्तिः। यथा पच्यादयः 
शक्तिशक्तिमदपेक्षा प्रापणीयजन्मानो न स्वयं शक्तयः। चितिः पुनः स्वयमेव शक्तिर्नार्त्थान्तरमपेक्षते। तेन 
नित्यावस्थितत्वञ्चितिशब्दस्य चिन्मात्राभिधायिनः शक्तिशब्देन सामानाधिकरण्यं, अविकृतरूपया एव 
चितेर्व्विषयित्वप्रदर्शनार्त्थम्।यस्मादेवमतो ऽपरिणामिनि व्यवस्थिते धर्म्मिणि धर्म्मान्तरतिरोभावेन धर्म्मान्तरप्रादुर्भावः 
परिणामः। न परिणन्तुं शीलमस्या इत्यपरिणामिनी॥

अत एवाप्रतिसंक्रमा परिणामिन एव, चित्तादेर्विषयादौ प्रतिसंक्रमदर्शनात्॥

अत एव दर्शितविषया दर्शितो ऽन्तःकरणेन विषयो ऽस्या इति तेनैव शुद्धा॥

अत एवानन्ता देशतः कालतश्च। पूर्व्वम्पूर्वमुत्तरस्योत्तरस्य हेतुत्वेन द्रष्टव्यम्। वैधर्म्म्यदृष्टान्तञ्चित्तमिन्द्रियाणि च। 
सत्वगुणातिमिका सत्वमेव गुणः सत्वगुणः तस्येवात्मा स्वरूपमवभासो ऽस्याः सेयं सत्वगुणात्मकावभासस्वरूपेत्यर्थः। 
वृत्यविशिष्टवृत्तित्वाद्वा सत्वगुणात्मिकेत्युच्यते॥

यदि वा ख्यात्या सम्बध्यते---ख्यातिः सत्वगुणात्मिकेति। अतश्चितिशक्तेरुक्तलक्षणाया विपरीता विलक्षणा विवेकख्यातिः 
परिणामादिमती॥

यस्मात्पुरुषादुत्कृष्टात्परिणामादिगुणविरहितान्निकृष्टा परिणामादिगुणयुक्तातः स्वरूपदोषदर्शनाद्विरक्तमपरक्तञ्चित्तन्तामपि 
ख्यातिमात्मनो निरुणाद्धि। तदवस्थन्निरोधावस्थं संस्कारोपगं संस्कारमात्रावशेषम्। निरुद्धासु वृत्तिषु वृत्तिजनिताः संस्कारा 
एवावशिष्यन्ते। एतस्यान्निरोधभूमौ यः समाधिः स निर्बीजः। निर्गतम्बीजमत्र \kle शादिबीजं सर्व्वमुत्सन्नमस्मिन्निति॥

तस्मादसंप्रज्ञातसमाधिलक्षणार्त्थमेवं सूत्रमित्युपसंहरति---स योगश्चित्तवृत्तिनिरोध इति॥२॥

% testing for Arlo

प्राङ्गणे

प्र\XeTeXglyph 204 णे

क्लेश

ङ्न